\chapter{Topología general}\label{ch:b3:topology}

El volumen de segundo año hizo análisis en espacios métricos:
distancias, bolas, sucesiones. Pero las nociones fundamentales --- la
\omterm{def:b3:topology:continuity}{continuidad}, la \omterm{def:b3:topology:compact}{compacidad}, la
\omterm{def:b3:topology:connected}{conexión} --- nunca mencionan el
valor numérico de una distancia, sino solo la familia de
\emph{\omterm{def:b3:topology:topology}{abiertos}} que esta genera. Este
capítulo toma esa familia como objeto primitivo. La ganancia no es
generalidad por sí misma: las construcciones cociente (el círculo como
$\R/\Z$, los espacios proyectivos), los productos y las \omterm{def:b3:topology:topology}{topologías} de
tipo débil del análisis funcional sencillamente no son objetos métricos
de nacimiento. Reconstruimos la
\omterm{def:b3:topology:continuity}{continuidad}, tratamos después la
\omterm{def:b3:topology:compact}{compacidad} mediante recubrimientos \omterm{def:b3:topology:topology}{abiertos} (demostrando que equivale a
la definición secuencial de segundo año en espacios métricos) y la
\omterm{def:b3:topology:connected}{conexión} --- para cerrar con el
teorema de que ninguna biyección
\omterm{def:b3:topology:continuity}{continua} puede identificar $\R$ con
$\R^2$: la \omterm{def:b3:topology:topology}{topología} sabe distinguir
dimensiones.

\section{Topologías, abiertos, continuidad}

\begin{definition}\label{def:b3:topology:topology}
Una \emph{topología}\index{topología} sobre un conjunto $X$ es una
familia $\mathcal T$ de subconjuntos de $X$ --- llamados
\emph{abiertos}\index{abierto} --- tal que: $\varnothing, X \in
\mathcal T$; toda unión de abiertos es abierta; toda intersección
\emph{finita} de abiertos es abierta. El par $(X, \mathcal T)$ es un
\emph{espacio topológico}. Los complementarios de los abiertos son los
\emph{cerrados}. Un \emph{entorno}\index{entorno} de $x$ es un conjunto
que contiene un abierto que contiene $x$.
\end{definition}

\begin{example}\label{ex:b3:topology:examples}
(a) Un espacio métrico, con «\omterm{def:b3:topology:topology}{abierto}» en el sentido de segundo año
(uniones de bolas abiertas): la
\emph{\omterm{def:b3:topology:topology}{topología} métrica}; métricas
distintas pueden dar la misma
\omterm{def:b3:topology:topology}{topología} (métricas equivalentes). Un
espacio cuya \omterm{def:b3:topology:topology}{topología} procede de
alguna métrica es \emph{metrizable}. (b) La
\omterm{def:b3:topology:topology}{topología} \emph{discreta} (todos los
subconjuntos) y la \omterm{def:b3:topology:topology}{topología}
\emph{indiscreta} $\{\varnothing, X\}$. (c) La
\omterm{def:b3:topology:topology}{topología} \emph{cofinita} sobre un
conjunto infinito: $=$ es \omterm{def:b3:topology:topology}{abierto} si es vacío o cofinito. No es
metrizable, como veremos
(\cref{exo:b3:topology:3}).
(d) Sobre $\R$, la \omterm{def:b3:topology:topology}{topología} usual;
sobre $\bar\R = \R \cup
\{\pm\infty\}$, la \omterm{def:b3:topology:topology}{topología} del
orden generada por las semirrectas --- lo que convierte
«$x_n \to +\infty$» en un caso de convergencia ordinaria.
\end{example}

\begin{definition}\label{def:b3:topology:interior}
Para $A \subseteq X$: el \emph{interior} $\mathring A$ es el mayor
\omterm{def:b3:topology:topology}{abierto} contenido en $A$ (la unión de
todos ellos); la \emph{clausura}\index{clausura} $\bar A$, el menor
cerrado que contiene $A$; la \emph{frontera} $\partial A = \bar A
\setminus \mathring A$. $A$ es \emph{denso}\index{conjunto denso} si
$\bar A = X$. Se tiene $x \in \bar A$ si y solo si todo
\omterm{def:b3:topology:topology}{entorno} de $x$ corta a $A$ (si algún
\omterm{def:b3:topology:topology}{entorno} no corta a $A$, el
complementario de su núcleo \omterm{def:b3:topology:topology}{abierto} es un cerrado más pequeño alrededor
de $A$; y recíprocamente).
\end{definition}

\begin{definition}\label{def:b3:topology:basis}
Una \emph{base}\index{base de una topología} de $\mathcal T$ es una
familia $\mathcal B \subseteq \mathcal T$ tal que todo
\omterm{def:b3:topology:topology}{abierto} es unión de miembros de
$\mathcal B$ (por ejemplo, las bolas abiertas en un espacio métrico o
los intervalos \omterm{def:b3:topology:topology}{abiertos} en $\R$). Una familia $\mathcal B$ de
subconjuntos de $X$ es base de \emph{alguna}
\omterm{def:b3:topology:topology}{topología} si y solo si recubre $X$ y,
para $B_1, B_2 \in \mathcal B$ y $x \in B_1 \cap B_2$, existe
$B_3 \in \mathcal B$ con $x \in B_3 \subseteq B_1\cap B_2$ --- y
entonces «las uniones de miembros» forman una
\omterm{def:b3:topology:topology}{topología}, la
\omterm{def:b3:topology:topology}{topología} \emph{generada} por
$\mathcal B$.
\end{definition}

\begin{definition}\label{def:b3:topology:continuity}
$f \colon X \to Y$ es \emph{continua}\index{aplicación continua} si
$f^{-1}(V)$ es \omterm{def:b3:topology:topology}{abierto} para todo \omterm{def:b3:topology:topology}{abierto} $V \subseteq Y$ ---
equivalentemente, las preimágenes de cerrados son cerradas;
equivalentemente, para todo $x$ y todo
\omterm{def:b3:topology:topology}{entorno} $V$ de $f(x)$, $f^{-1}(V)$ es
un \omterm{def:b3:topology:topology}{entorno} de $x$ (continuidad
\emph{en} cada $x$). Basta comprobarlo sobre una
\omterm{def:b3:topology:basis}{base} de $Y$. Las composiciones de
aplicaciones continuas son continuas. Un
\emph{homeomorfismo}\index{homeomorfismo} es una biyección continua con
inversa continua; la \omterm{def:b3:topology:topology}{topología}
estudia las propiedades que los homeomorfismos conservan.
\end{definition}

\begin{proposition}[Continuidad y clausura; pegado]
\label{prop:b3:topology:gluing}
(a) $f$ es \omterm{def:b3:topology:continuity}{continua} si y solo si
$f(\bar A) \subseteq \overline{f(A)}$ para todo $A \subseteq X$. (b) Si
$X = F_1 \cup F_2$ con $F_1, F_2$ cerrados y $f\colon X
\to Y$ se restringe \omterm{def:b3:topology:continuity}{continuamente} a
cada $F_i$, entonces $f$ es
\omterm{def:b3:topology:continuity}{continua}.
\end{proposition}

\begin{proof}
(a) Si $f$ es \omterm{def:b3:topology:continuity}{continua}:
$f^{-1}(\overline{f(A)})$ es cerrado y contiene $A$, luego contiene
$\bar A$. Recíprocamente, aplíquese el criterio a $A = f^{-1}(F)$ con
$F$ cerrado: $f(\bar A) \subseteq
\overline{f(f^{-1}(F))} \subseteq \bar F = F$, luego $\bar A
\subseteq f^{-1}(F) = A$: las preimágenes de cerrados son cerradas. (b)
Para $F \subseteq Y$ cerrado: $f^{-1}(F) =
(f\restriction_{F_1})^{-1}(F) \cup (f\restriction_{F_2})^{-1}
(F)$, unión de dos conjuntos cerrados en $F_1$ y en $F_2$
respectivamente, luego cerrada en $X$ ($F_i$ cerrados: cerrado dentro de
cerrado es cerrado).
\end{proof}

\begin{definition}\label{def:b3:topology:hausdorff}
$X$ es \emph{de Hausdorff}\index{espacio de Hausdorff} (o
\emph{separado}) si dos puntos distintos cualesquiera tienen
\omterm{def:b3:topology:topology}{entornos} disjuntos. Los espacios
métricos son de Hausdorff (bolas de radio $d(x,y)/2$). Una sucesión
$(x_n)$ \emph{converge} a $x$ si todo
\omterm{def:b3:topology:topology}{entorno} de $x$ contiene todos los
$x_n$ salvo un número finito; en un espacio de Hausdorff el límite es
único (dos límites tendrían \omterm{def:b3:topology:topology}{entornos}
disjuntos, cada uno con una cola de la sucesión). En un espacio de
Hausdorff, los puntos --- y por tanto los conjuntos finitos --- son
cerrados.
\end{definition}

\begin{remark}\label{rem:b3:topology:sequences}
En los espacios métricos, las sucesiones lo detectan todo:
$x \in \bar A$ si y solo si alguna sucesión de $A$ converge a $x$
(tómese $x_n \in A \cap
B(x, 1/n)$), y $f$ es \omterm{def:b3:topology:continuity}{continua} si y
solo si es secuencialmente
\omterm{def:b3:topology:continuity}{continua}. En espacios generales
fallan ambas equivalencias; los sustitutos correctos de las sucesiones
(filtros, redes) corresponden a un curso más avanzado. Enunciaremos los
resultados basados en sucesiones en espacios métricos y los basados en
recubrimientos en general --- y demostraremos su equivalencia allí donde
la hay.
\end{remark}

\section{Subespacios, productos, cocientes}

\begin{definition}\label{def:b3:topology:constructions}
Tres maneras de fabricar espacios nuevos a partir de
otros:\index{topología producto}\index{topología cociente}
\begin{enumerate}
  \item \emph{Subespacio}: sobre $A \subseteq X$, los
  \omterm{def:b3:topology:topology}{abiertos} son los $U \cap A$ con $U$
  \omterm{def:b3:topology:topology}{abierto} en $X$ --- la \omterm{def:b3:topology:topology}{topología}
  menos fina que hace \omterm{def:b3:topology:continuity}{continua} la
  inclusión. \item \emph{Producto}: sobre $X \times Y$ (y sobre los
  productos finitos), la \omterm{def:b3:topology:topology}{topología}
  cuya \omterm{def:b3:topology:basis}{base} son las \emph{cajas
  abiertas} $U \times
        V$; en un producto infinito $\prod_i X_i$, la
  \omterm{def:b3:topology:basis}{base} está formada por las cajas
  $\prod U_i$ con $U_i = X_i$ para \emph{todos los $i$ salvo un número
  finito} --- la \omterm{def:b3:topology:topology}{topología} menos fina
  que hace \omterm{def:b3:topology:continuity}{continuas} todas las
  proyecciones. \item \emph{Cociente}: si $\sim$ es una equivalencia
  sobre $X$ y $\pi\colon X \to X/{\sim}$ la proyección, declárese $V
        \subseteq X/{\sim}$ \omterm{def:b3:topology:topology}{abierto} si y solo si $\pi^{-1}(V)$ es
  \omterm{def:b3:topology:topology}{abierto} --- la \omterm{def:b3:topology:topology}{topología} más fina
  que hace $\pi$ \omterm{def:b3:topology:continuity}{continua}.
\end{enumerate}
\end{definition}

\begin{proposition}[Propiedades universales]
\label{prop:b3:topology:universal}
(a) $f \colon Z \to X \times Y$ es
\omterm{def:b3:topology:continuity}{continua} si y solo si lo son sus
dos componentes $\pi_X \circ f$, $\pi_Y \circ f$ (lo mismo para
productos arbitrarios). (b) $g \colon X/{\sim} \to Z$ es
\omterm{def:b3:topology:continuity}{continua} si y solo si lo es
$g \circ \pi
\colon X \to Z$.
\end{proposition}

\begin{proof}
(a) Necesidad: composiciones. Suficiencia: basta comprobar las
preimágenes de las cajas de la \omterm{def:b3:topology:basis}{base}:
$f^{-1}(U \times V) = (\pi_Xf)^{-1}(U)
\cap (\pi_Yf)^{-1}(V)$, abierta (solo un número finito de factores
$\neq X_i$ en el caso infinito). (b) Necesidad: composición.
Suficiencia: para $W \subseteq Z$ \omterm{def:b3:topology:topology}{abierto},
$\pi^{-1}(g^{-1}(W)) = (g\pi)^{-1}(W)$ es \omterm{def:b3:topology:topology}{abierto}, lo que por definición
de la \omterm{def:b3:topology:constructions}{topología cociente}
significa que $g^{-1}(W)$ es \omterm{def:b3:topology:topology}{abierto}.
\end{proof}

\begin{example}\label{ex:b3:topology:circle}
El cociente $\R/\Z$ (identificando $x$ con $x + n$) es
\omterm{def:b3:topology:continuity}{homeomorfo} al círculo
$S^1 = \{z \in \C : \abs z = 1\}$: la aplicación $x
\mapsto \eu^{2\iu\pi x}$ pasa al cociente y da una biyección
\omterm{def:b3:topology:continuity}{continua} $\R/\Z \to S^1$
(la~\cref{prop:b3:topology:universal}(b)); su inversa es
\omterm{def:b3:topology:continuity}{continua} por el argumento de
\omterm{def:b3:topology:compact}{compacidad} del~\cref{cor:b3:topology:compacthomeo} de más abajo
(el~\cref{exo:b3:topology:5} lo detalla todo, incluido por qué $\R/\Z$
es \omterm{def:b3:topology:hausdorff}{de Hausdorff} y
\omterm{def:b3:topology:compact}{compacto}). Del mismo modo, $[0,1]$ con
los extremos pegados es $S^1$, el cuadrado con los lados opuestos
pegados es el toro, y el pegado pasa por fin a ser un teorema, y no un
dibujo.
\end{example}

\section{Compacidad}

\begin{definition}\label{def:b3:topology:compact}
Un \emph{recubrimiento \omterm{def:b3:topology:topology}{abierto}} de $X$ es una familia $(U_i)_{i\in I}$
de \omterm{def:b3:topology:topology}{abiertos} con $\bigcup U_i = X$.
$X$ es \emph{compacto}\index{espacio compacto} si es
\omterm{def:b3:topology:hausdorff}{de Hausdorff} y todo recubrimiento
\omterm{def:b3:topology:topology}{abierto} admite un \emph{subrecubrimiento finito}. Equivalentemente
(tomando complementarios): toda familia de cerrados con la
\emph{propiedad de intersección finita} (toda subfamilia finita tiene
intersección no vacía) tiene intersección total no vacía.
\end{definition}

\begin{theorem}[Primeras
propiedades]\label{thm:b3:topology:compactprops}
Sea $X$ \omterm{def:b3:topology:compact}{compacto}.
\begin{enumerate}
  \item Un subconjunto cerrado de $X$ es
  \omterm{def:b3:topology:compact}{compacto}; un subconjunto
  \omterm{def:b3:topology:compact}{compacto} de un
  \omterm{def:b3:topology:hausdorff}{espacio de Hausdorff} es cerrado.
  \item La imagen \omterm{def:b3:topology:continuity}{continua} de un
  \omterm{def:b3:topology:compact}{espacio compacto} dentro de un
  \omterm{def:b3:topology:hausdorff}{espacio de Hausdorff} es
  compacta. En particular, una
  $f\colon X
        \to \R$ \omterm{def:b3:topology:continuity}{continua} está
  acotada y alcanza sus cotas. \item Una sucesión decreciente de
  cerrados no vacíos de $X$ tiene intersección no vacía.
\end{enumerate}
\end{theorem}

\begin{proof}
(1) Sea $F \subseteq X$ cerrado y sea $(U_i)$ un recubrimiento \omterm{def:b3:topology:topology}{abierto}
de $F$ (por \omterm{def:b3:topology:topology}{abiertos} de $X$): añadiendo $X \setminus F$ se obtiene un
recubrimiento \omterm{def:b3:topology:topology}{abierto} de $X$; un subrecubrimiento finito, quitando
$X \setminus F$, recubre $F$. Que el subespacio sea
\omterm{def:b3:topology:hausdorff}{de Hausdorff} es claro.
Recíprocamente, sean $K \subseteq Y$
\omterm{def:b3:topology:compact}{compacto}, $Y$
\omterm{def:b3:topology:hausdorff}{de Hausdorff} y $y \notin K$: para
cada $x \in K$ tómense \omterm{def:b3:topology:topology}{abiertos} disjuntos $U_x \ni x$, $V_x \ni y$; un
número finito de $U_x$ recubre $K$, y la intersección de los $V_x$
correspondientes es un \omterm{def:b3:topology:topology}{entorno} de $y$
disjunto del recubrimiento de $K$ y, por tanto, de $K$: el
complementario de $K$ es \omterm{def:b3:topology:topology}{abierto}.

(2) Si $(V_j)$ recubre $f(X)$, entonces $(f^{-1}(V_j))$ recubre $X$; un
subrecubrimiento finito abajo procede del subrecubrimiento finito de
arriba. La imagen es \omterm{def:b3:topology:hausdorff}{de Hausdorff}
por ser subespacio. Para $f$ real: $f(X)$ es
\omterm{def:b3:topology:compact}{compacto} en $\R$, luego cerrado y
acotado (recúbrase por $(-n, n)$ para la acotación; cerrado por (1)), y
un conjunto cerrado y acotado contiene su supremo.

(3) Si $\bigcap F_n = \varnothing$, los \omterm{def:b3:topology:topology}{abiertos} $X \setminus F_n$
recubren $X$; bastan un número finito, de modo que algún $F_{n_0} =
\varnothing$ (la sucesión decrece) --- contradicción.
\end{proof}

\begin{corollary}\label{cor:b3:topology:compacthomeo}
Una \emph{biyección} \omterm{def:b3:topology:continuity}{continua} de un
\omterm{def:b3:topology:compact}{espacio compacto} en un
\omterm{def:b3:topology:hausdorff}{espacio de Hausdorff} es un
\omterm{def:b3:topology:continuity}{homeomorfismo}.
\end{corollary}

\begin{proof}
La inversa es \omterm{def:b3:topology:continuity}{continua} si y solo si
las imágenes directas de cerrados son cerradas; un $F$ cerrado es
\omterm{def:b3:topology:compact}{compacto}
(el~\cref{thm:b3:topology:compactprops}(1)), su imagen es
compacta (2) y, por tanto, cerrada (1)
en el destino \omterm{def:b3:topology:hausdorff}{de Hausdorff}.
\end{proof}

\begin{theorem}[Productos finitos]\label{thm:b3:topology:tychonoff}
Un producto finito de \omterm{def:b3:topology:compact}{espacios
compactos} es \omterm{def:b3:topology:compact}{compacto}.\index{teorema
de Tychonoff (caso finito)}
\end{theorem}

\begin{proof}
Basta tratar $X \times Y$. Ser \omterm{def:b3:topology:hausdorff}{de
Hausdorff} se hereda (sepárese en una coordenada). Sea $(W_i)$ un
recubrimiento \omterm{def:b3:topology:topology}{abierto} de $X
\times Y$; podemos suponer que los $W_i$ son cajas de la
\omterm{def:b3:topology:basis}{base} $U_i \times
V_i$ (refínese: cada punto está en una caja contenida en algún $W_i$; un
subrecubrimiento finito por cajas da uno por $W_i$). Fijemos $x \in X$:
la \emph{rebanada} $\{x\}\times Y \cong Y$ es
compacta, luego un número finito de
cajas $U_1\times V_1, \dots, U_k\times V_k$ la recubre, con $x \in U_j$
para todo $j$; entonces $U^x = \bigcap_{j \leq k} U_j$ es un
\omterm{def:b3:topology:topology}{entorno} \omterm{def:b3:topology:topology}{abierto} de $x$ con
$U^x \times Y$ recubierto por un número finito de cajas (el \emph{lema
del tubo}\index{lema del tubo}: para $(x', y) \in U^x \times Y$,
$y \in V_j$ para algún $j$, pues las cajas recubrían $\{x\}\times Y$ al
nivel $y$, y $x' \in U^x \subseteq U_j$, luego $(x', y) \in U_j
\times V_j$). Ahora bien, un número finito de $U^x$ recubre el
\omterm{def:b3:topology:compact}{compacto} $X$; las colecciones finitas
de cajas correspondientes recubren $X \times Y$.
\end{proof}

\begin{theorem}[Compacidad en espacios métricos]
\label{thm:b3:topology:metriccompact}
Para un espacio métrico $(X, d)$, son equivalentes:\index{compacidad
secuencial}
\begin{enumerate}
  \item $X$ es \omterm{def:b3:topology:compact}{compacto}
  (Borel--Lebesgue); \item toda sucesión de $X$ tiene una subsucesión
  convergente (\omterm{thm:b3:topology:metriccompact}{compacidad
  secuencial} --- la definición de segundo año); \item $X$ es completo y
  \emph{totalmente acotado}: para todo $\varepsilon > 0$, un número
  finito de bolas de radio $\varepsilon$ recubre $X$.
\end{enumerate}
\end{theorem}

\begin{proof}
(1)$\Rightarrow$(2): sea $(x_n)$ sin subsucesión convergente. Entonces
cada $x \in X$ tiene una bola abierta $B_x$ que contiene $x_n$ solo para
un número finito de $n$ (si no, alguna subsucesión convergería a $x$:
tómense radios $1/k$). Un número finito de $B_x$ recubre $X$, luego solo
existe un número finito de índices $n$: absurdo.

(2)$\Rightarrow$(3): completitud: una sucesión de Cauchy con una
subsucesión convergente converge (segundo año). Acotación total: si
algún $\varepsilon$ no admite recubrimiento finito, elíjanse
inductivamente $x_{n+1}$ fuera de
$\bigcup_{k\leq n} B(x_k, \varepsilon)$: la sucesión cumple
$d(x_m, x_n) \geq \varepsilon$ para $m \neq n$, no tiene subsucesión de
Cauchy ni, por tanto, convergente.

(3)$\Rightarrow$(1): primero, (3) implica (2): dado $(x_n)$, recúbrase
$X$ por un número finito de bolas de radio $1$: una de ellas, $B_1$,
contiene una subsucesión; recúbrase $X$ por bolas de radio $1/2$: una
contiene una subsucesión ulterior; itérese y diagonalícese: la
subsucesión diagonal es de Cauchy (dos términos posteriores a la etapa
$k$ están en una misma bola de radio $2^{-k}$, salvo el $2\cdot$
habitual) y, por tanto, converge. Sea ahora $(U_i)$ un recubrimiento
\omterm{def:b3:topology:topology}{abierto} y supóngase que no hay subrecubrimiento finito. \emph{Argumento
del número de Lebesgue}: para cada $n$, alguna bola $B(y_n, 2^{-n})$ no
queda recubierta por un número finito de $U_i$ --- en efecto, recúbrase
$X$ por un número finito de bolas de radio $2^{-n}$; si cada una
admitiera recubrimiento finito, también lo admitiría $X$. Por (2), una
subsucesión $y_{n_k} \to y$; tómense $i$ con $y \in U_i$ y $r > 0$ con
$B(y, r) \subseteq U_i$. Para $k$ grande,
$B(y_{n_k}, 2^{-n_k}) \subseteq B(y, r) \subseteq U_i$: queda recubierta
por \emph{un} solo $U_i$ --- contradicción.
\end{proof}

\begin{corollary}[Heine--Borel; Heine]
\label{cor:b3:topology:heineborel}
(a) Un subconjunto de $\R^n$ es
\omterm{def:b3:topology:compact}{compacto} si y solo si es cerrado y
acotado. (b) Una \omterm{def:b3:topology:continuity}{aplicación
continua} de un espacio métrico
\omterm{def:b3:topology:compact}{compacto} en un espacio métrico es
uniformemente
\omterm{def:b3:topology:continuity}{continua}.\index{teorema de
Heine--Borel}\index{teorema de Heine}
\end{corollary}

\begin{proof}
(a) Cerrado y acotado $\Rightarrow$ contenido en un cubo $[-M,M]^n$, que
es \omterm{def:b3:topology:compact}{compacto}: $[-M, M]$ lo es
(secuencialmente, por Bolzano--Weierstrass --- o directamente por
dicotomía para recubrimientos), y el~\cref{thm:b3:topology:tychonoff} se
ocupa del producto; aplíquese después
el~\cref{thm:b3:topology:compactprops}(1). Recíprocamente, un
subconjunto \omterm{def:b3:topology:compact}{compacto} es cerrado
(el~\cref{thm:b3:topology:compactprops}(1)) y acotado (recúbrase por
bolas concéntricas).

(b) Sean $f \colon X \to Y$, $\varepsilon > 0$. Las bolas
$B\bigl(x, \delta_x\bigr)$ con $f\bigl(B(x,
2\delta_x)\bigr)\subseteq B\bigl(f(x), \varepsilon/2\bigr)$ recubren
$X$; extráigase un subrecubrimiento finito $B(x_j, \delta_{x_j})$ y
póngase $\delta = \min_j \delta_{x_j}$. Si $d(x, x') < \delta$: $x
\in B(x_j, \delta_{x_j})$ para algún $j$ y entonces ambos $x, x'
\in B(x_j, 2\delta_{x_j})$, luego $d(f(x), f(x')) < \varepsilon$.
\end{proof}

\begin{definition}\label{def:b3:topology:locallycompact}
$X$ es \emph{localmente compacto}\index{espacio localmente compacto} si
es \omterm{def:b3:topology:hausdorff}{de Hausdorff} y todo punto tiene
un \omterm{def:b3:topology:topology}{entorno}
\omterm{def:b3:topology:compact}{compacto} ($\R^n$; los \omterm{def:b3:topology:topology}{abiertos} de
$\R^n$; los espacios discretos --- pero no $\Q$, véase
el~\cref{exo:b3:topology:9}). Todo
espacio localmente \omterm{def:b3:topology:compact}{compacto} se sumerge
en uno \omterm{def:b3:topology:compact}{compacto}: la
\emph{compactificación por un punto}\index{compactificación por un
punto} $\hat X = X
\cup \{\infty\}$, cuyos \omterm{def:b3:topology:topology}{abiertos} son los de $X$ junto con los
complementarios (en $\hat X$) de los subconjuntos
\omterm{def:b3:topology:compact}{compactos} de $X$. Los axiomas se
comprueban directamente; $\hat X$ es
\omterm{def:b3:topology:compact}{compacto} (un recubrimiento tiene un
miembro que contiene $\infty$, cuyo complementario es
\omterm{def:b3:topology:compact}{compacto} y queda recubierto por un
número finito de los demás) y \omterm{def:b3:topology:hausdorff}{de
Hausdorff} (sepárese $x$ de $\infty$ mediante un
\omterm{def:b3:topology:topology}{entorno}
\omterm{def:b3:topology:compact}{compacto} de $x$ y su complementario).
Ejemplo: $\hat\R \cong S^1$, y $\widehat{\R^n} \cong S^n$ por proyección
estereográfica (el~\cref{exo:b3:topology:11}).
\end{definition}

\section{Conexión}

\begin{definition}\label{def:b3:topology:connected}
$X$ es \emph{conexo}\index{espacio conexo} si no es unión de dos
\omterm{def:b3:topology:topology}{abiertos} disjuntos no vacíos ---
equivalentemente, si sus únicos subconjuntos a la vez \omterm{def:b3:topology:topology}{abiertos} y
cerrados son $\varnothing$ y $X$; equivalentemente, si toda
\omterm{def:b3:topology:continuity}{aplicación continua}
$X \to \{0, 1\}$ (con la \omterm{def:b3:topology:topology}{topología} discreta) es constante. Un
subconjunto es conexo si lo es como subespacio.
\end{definition}

\begin{theorem}\label{thm:b3:topology:connectedbasics}
\begin{enumerate}
  \item Los subconjuntos \omterm{def:b3:topology:connected}{conexos} de
  $\R$ son exactamente los intervalos. \item Las imágenes
  \omterm{def:b3:topology:continuity}{continuas} de
  \omterm{def:b3:topology:connected}{conexos} son
  \omterm{def:b3:topology:connected}{conexas} (de donde el teorema del
  valor intermedio: una aplicación real
  \omterm{def:b3:topology:continuity}{continua} sobre un
  \omterm{def:b3:topology:connected}{espacio conexo} tiene por imagen un
  intervalo). \item Si $(A_i)$ son
  \omterm{def:b3:topology:connected}{conexos} con un punto común,
  $\bigcup
        A_i$ es \omterm{def:b3:topology:connected}{conexo}. Si $A$ es
  \omterm{def:b3:topology:connected}{conexo} y $A \subseteq
        B \subseteq \bar A$, entonces $B$ es
  \omterm{def:b3:topology:connected}{conexo}. \item Los productos
  finitos de \omterm{def:b3:topology:connected}{espacios conexos} son
  \omterm{def:b3:topology:connected}{conexos}.
\end{enumerate}
\end{theorem}

\begin{proof}
En todo lo que sigue usamos el criterio de $\{0,1\}$: toda
$f \colon X \to
\{0,1\}$ \omterm{def:b3:topology:continuity}{continua} ha de ser
constante.

(1) Un $A$ que no sea un intervalo omite algún $c$ comprendido entre
$a, b \in A$:
$A = (A \cap (-\infty, c)) \sqcup (A \cap (c, +\infty))$
lo desconecta. Recíprocamente, sea $I$ un intervalo y sea $f \colon I
\to \{0,1\}$ \omterm{def:b3:topology:continuity}{continua} con
$f(a) = 0$, $f(b) = 1$, $a < b$. Sea
$c = \sup\{x \in [a,b] : f(x) = 0\}$; la
\omterm{def:b3:topology:continuity}{continuidad} en $c$ obliga a
$f(c) = 0$ (límite de valores $0$: todo
\omterm{def:b3:topology:topology}{entorno} de $c$ corta a $\{f = 0\}$, y
$\{f=0\}$ es cerrado) y entonces $c < b$ con $f \equiv 1$ sobre
$(c, b]$, luego $f(c) = 1$ por el mismo argumento de
\omterm{def:b3:topology:interior}{clausura} sobre $\{f = 1\} \ni c$:
contradicción.

(2) Una $g \colon f(X) \to \{0,1\}$
\omterm{def:b3:topology:continuity}{continua} da $g \circ f$ constante,
de modo que $g$ es constante sobre la imagen.

(3) Una $f\colon \bigcup A_i \to \{0,1\}$
\omterm{def:b3:topology:continuity}{continua} es constante en cada
$A_i$, con el mismo valor en el punto común. Para la
\omterm{def:b3:topology:interior}{clausura}: $f \colon B \to \{0,1\}$
vale constantemente $= c$ sobre $A$; todo $x \in B \subseteq \bar A$
está en la \omterm{def:b3:topology:interior}{clausura} de $A$, y $f^{-1}
(f(x))$ es un \omterm{def:b3:topology:topology}{entorno} de $x$
(preimagen de un \omterm{def:b3:topology:topology}{abierto}), que ha de cortar a $A$: $f(x) = c$.

(4) Para $X \times Y$ y $f\colon X\times Y \to \{0,1\}$: dos puntos
cualesquiera $(x,y), (x',y')$ se unen mediante el «codo»
$\{x\}\times Y \cup X \times \{y'\}$, unión de dos
\omterm{def:b3:topology:connected}{conexos}
(\omterm{def:b3:topology:continuity}{homeomorfos} a $Y$ y a $X$) que se
cortan en $(x, y')$: por (3) y (2), $f$ toma el mismo valor en los dos
puntos.
\end{proof}

\begin{definition}\label{def:b3:topology:pathconnected}
$X$ es \emph{conexo por caminos}\index{espacio conexo por caminos} si
dos puntos cualesquiera se unen por un \emph{camino} (una $\gamma
\colon [0,1] \to X$ \omterm{def:b3:topology:continuity}{continua}).
\omterm{def:b3:topology:connected}{Conexo} por caminos implica
\omterm{def:b3:topology:connected}{conexo}: dos valores de una
$f \colon X \to \{0,1\}$ \omterm{def:b3:topology:continuity}{continua}
en $x, y$ son valores de la constante $f\circ\gamma$
(el~\cref{thm:b3:topology:connectedbasics}(1)--(2)). Los subconjuntos
convexos de un espacio normado son
\omterm{def:b3:topology:connected}{conexos} por caminos (segmentos);
también lo son $\R^n
\setminus \{0\}$ para $n \geq 2$ (rodéese el origen) y $S^n$ para
$n \geq 1$ (proyéctense caminos desde $\R^{n+1}\setminus
\{0\}$).
\end{definition}

\begin{example}[La curva seno del topólogo]
\label{ex:b3:topology:sinecurve}
Sean $\Gamma = \{(x, \sin\frac1x) : 0 < x \leq 1\}$ y $S =
\bar\Gamma = \Gamma \cup (\{0\}\times[-1,1])$ (todo punto $(0,
y)$ con $\abs y \leq 1$ es límite de puntos de $\Gamma$: resuélvase
$\sin\frac1x = y$ cerca de $0$). Entonces $S$ es
\emph{\omterm{def:b3:topology:connected}{conexo}} ---
\omterm{def:b3:topology:interior}{clausura} del
\omterm{def:b3:topology:connected}{conexo} $\Gamma$, imagen
\omterm{def:b3:topology:continuity}{continua} de $(0, 1]$
(el~\cref{thm:b3:topology:connectedbasics}(3)) --- pero \emph{no es
\omterm{def:b3:topology:pathconnected}{conexo por caminos}}: un
\omterm{def:b3:topology:pathconnected}{camino} de $(1, \sin 1)$ a
$(0,0)$ tendría que recorrer abscisas $\to 0$ mientras la ordenada
oscila entre $\pm1$; el~\cref{exo:b3:topology:9} lo hace riguroso. La
\omterm{def:b3:topology:connected}{conexión} y la
\omterm{def:b3:topology:connected}{conexión} por \omterm{def:b3:topology:pathconnected}{caminos} son
genuinamente distintas.\index{curva seno del topólogo}
\end{example}

\begin{omfigure}
\begin{tikzpicture}
\begin{axis}[omaxis, width=10.5cm, height=4.6cm,
    xmin=-0.02, xmax=0.55, ymin=-1.25, ymax=1.25,
    xtick={0, 0.25, 0.5}, ytick={-1, 1}]
  \addplot[omThm, thick, domain=0.006:0.55, samples=1200]
    {sin(deg(1/x))};
  \addplot[omDef, very thick] coordinates {(0,-1) (0,1)};
\end{axis}
\end{tikzpicture}

{\small La \omterm{ex:b3:topology:sinecurve}{curva seno del topólogo}:
la gráfica de $\sin\frac1x$ (en rojo) se acumula sobre todo el segmento
$\{0\}\times[-1,1]$ (en azul). La unión es
\omterm{def:b3:topology:connected}{conexa} pero no
\omterm{def:b3:topology:pathconnected}{conexa por caminos}: ningún
\omterm{def:b3:topology:pathconnected}{camino}
\omterm{def:b3:topology:continuity}{continuo} puede atravesar las
infinitas oscilaciones en un tiempo de parámetro finito.}
\end{omfigure}

\begin{definition}\label{def:b3:topology:components}
La \emph{componente conexa}\index{componente conexa} de $x
\in X$ es la unión de todos los subconjuntos
\omterm{def:b3:topology:connected}{conexos} que contienen $x$ --- el
mayor de ellos (el~\cref{thm:b3:topology:connectedbasics}(3)). Las
componentes forman una partición de $X$ y son cerradas (las
\omterm{def:b3:topology:interior}{clausuras} de
\omterm{def:b3:topology:connected}{conexos} son
\omterm{def:b3:topology:connected}{conexas}). $X$ es \emph{totalmente
disconexo} si todas sus componentes se reducen a un punto ($\Q$; el
conjunto de Cantor del problema de fin de semana).
\end{definition}

\begin{theorem}\label{thm:b3:topology:rnotr2}
$\R$ no es \omterm{def:b3:topology:continuity}{homeomorfo} a ningún
$\R^n$ con $n \geq 2$.
\end{theorem}

\begin{proof}
Supongamos que $h \colon \R \to \R^n$ es un
\omterm{def:b3:topology:continuity}{homeomorfismo}. Al suprimir un
punto: $\R \setminus \{0\}$ es
\omterm{def:b3:topology:continuity}{homeomorfo} a $\R^n \setminus
\{h(0)\}$. Pero $\R\setminus\{0\}$ es disconexo, mientras que
$\R^n\setminus\{\text{punto}\}$ es
\omterm{def:b3:topology:pathconnected}{conexo por caminos} para $n \geq
2$ (la~\cref{def:b3:topology:pathconnected}; trasládese el punto a $0$):
la \omterm{def:b3:topology:connected}{conexión} es un invariante por
\omterm{def:b3:topology:continuity}{homeomorfismos}
(el~\cref{thm:b3:topology:connectedbasics}(2)) --- contradicción. (El
caso $\R^2 \not\cong \R^3$ exige invariantes más finos --- la
\omterm{def:b3:topology:topology}{topología}
\omterm{def:b3:galois:algebraic}{algebraica}; el problema de fin de
semana muestra el peligro: sí existen \emph{sobreyecciones
\omterm{def:b3:topology:continuity}{continuas}} $\R \to \R^2$.)
\end{proof}

\begin{method}\label{met:b3:topology:connectargument}
Para demostrar que un conjunto es
\omterm{def:b3:topology:connected}{conexo}: exhíbase como imagen
\omterm{def:b3:topology:continuity}{continua}, como unión de
\omterm{def:b3:topology:connected}{conexos} que se solapan, como
\omterm{def:b3:topology:interior}{clausura} o como producto
(el~\cref{thm:b3:topology:connectedbasics}); para subconjuntos de
espacios normados, demuéstrese la
\omterm{def:b3:topology:connected}{conexión} por \omterm{def:b3:topology:pathconnected}{caminos} con
\omterm{def:b3:topology:pathconnected}{caminos} explícitos (segmentos,
arcos, codos). Para demostrar que dos espacios \emph{no} son
\omterm{def:b3:topology:continuity}{homeomorfos}: búsquese un invariante
\omterm{def:b3:topology:topology}{topológico} que difiera ---
\omterm{def:b3:topology:compact}{compacidad}, \omterm{def:b3:topology:connected}{conexión}, número de
componentes o componentes tras suprimir un conjunto finito bien elegido
(el truco del~\cref{thm:b3:topology:rnotr2}: también demuestra que
$[0,1) \not\cong
(0,1)$ y $S^1 \not\cong [0,1]$).
\end{method}

\section{Ejercicios}

\begin{exercise}[$\star$]\label{exo:b3:topology:1}
(a) Listar todas las \omterm{def:b3:topology:topology}{topologías} sobre $\{a, b\}$, clasificarlas salvo
\omterm{def:b3:topology:continuity}{homeomorfismo} y determinar cuáles
son \omterm{def:b3:topology:connected}{conexas} y cuáles
\omterm{def:b3:topology:hausdorff}{de Hausdorff}. (b) En $\R$ (con la
\omterm{def:b3:topology:topology}{topología} usual), calcular el
\omterm{def:b3:topology:interior}{interior}, la
\omterm{def:b3:topology:interior}{clausura} y la frontera de $\Q$, de
$[0,1) \cup \{2\}$ y de $\{1/n : n \geq
1\}$.
\end{exercise}

\begin{exercise}[$\star$]\label{exo:b3:topology:2}
(a) Demostrar que $f\colon X \to Y$ es
\omterm{def:b3:topology:continuity}{continua} si y solo si $f^{-1}(B)$
es \omterm{def:b3:topology:topology}{abierto} para todo $B$ de una \omterm{def:b3:topology:basis}{base}
fija de $Y$. (b) Demostrar que
$f = \mathbf 1_{[0,\infty)} \colon \R \to \R$ no es
\omterm{def:b3:topology:continuity}{continua} pero sí es
\emph{\omterm{def:b3:topology:continuity}{continua} por la derecha}.
Comprobar que los intervalos semiabiertos $[a, b)$ forman una
\omterm{def:b3:topology:basis}{base de una topología} sobre el origen
$\R$ (la \emph{recta de Sorgenfrey}), que esta
\omterm{def:b3:topology:topology}{topología} es estrictamente más fina
que la usual, y que una función $\R \to
\R$ (con el destino usual) es
\omterm{def:b3:topology:continuity}{continua} desde la recta de
Sorgenfrey si y solo si es \omterm{def:b3:topology:continuity}{continua}
por la derecha en todo punto.
\end{exercise}

\begin{exercise}[$\star$]\label{exo:b3:topology:3}
Sobre un conjunto infinito con la
\omterm{def:b3:topology:topology}{topología} cofinita, demostrar: dos
\omterm{def:b3:topology:topology}{abiertos} no vacíos cualesquiera se
cortan (de modo que el espacio no es
\omterm{def:b3:topology:hausdorff}{de Hausdorff} y, por tanto, no es
metrizable); toda sucesión inyectiva converge a \emph{todos} los puntos.
¿Dónde usa la demostración de la unicidad del límite la propiedad
\omterm{def:b3:topology:hausdorff}{de Hausdorff}?
\end{exercise}

\begin{exercise}[$\star\star$]\label{exo:b3:topology:4}
(a) Demostrar que las proyecciones $\pi_X, \pi_Y$ de un producto son
\omterm{def:b3:topology:continuity}{continuas} y \emph{abiertas} (las
imágenes de \omterm{def:b3:topology:topology}{abiertos} son abiertas), pero no cerradas en general
($\{xy = 1\}$ en $\R^2$). (b) Demostrar que una sucesión de un producto
numerable $\prod_n X_n$ de espacios métricos converge si y solo si
converge cada coordenada, y que
$d(x, y) = \sum_n 2^{-n}\min(1, d_n(x_n, y_n))$ metriza la
\omterm{def:b3:topology:constructions}{topología producto}. (c) Sobre
$\R^\N$, demostrar que la «\omterm{def:b3:topology:topology}{topología}
de cajas» (todo producto de \omterm{def:b3:topology:topology}{abiertos} es \omterm{def:b3:topology:topology}{abierto}) es estrictamente más
fina: la sucesión $x^{(k)} = (1/k, 1/k, \dots)$ converge a $0$ en la
\omterm{def:b3:topology:constructions}{topología producto} pero no en la
\omterm{def:b3:topology:topology}{topología} de cajas.
\end{exercise}

\begin{exercise}[$\star\star$]\label{exo:b3:topology:5}
El círculo, de tres maneras. Demostrar que los siguientes espacios son
\omterm{def:b3:topology:continuity}{homeomorfos} dos a dos, con
aplicaciones explícitas: (i) $S^1 \subseteq \R^2$; (ii) $\R/\Z$
(\omterm{def:b3:topology:constructions}{topología cociente}); (iii)
$[0,1]/(0 \sim 1)$. \emph{(Para (ii): demuéstrese que $\R/\Z$ es
\omterm{def:b3:topology:hausdorff}{de Hausdorff} --- levántense dos
clases a representantes a distancia $\leq \frac12$ --- y que
$\pi([0,1])$ es todo el espacio, de modo que el cociente es
\omterm{def:b3:topology:compact}{compacto}; después úsese
\cref{cor:b3:topology:compacthomeo}.)}
\end{exercise}

\begin{exercise}[$\star\star$]\label{exo:b3:topology:6}
Sea $X$ un espacio métrico. (a) Demostrar que una unión finita de
\omterm{def:b3:topology:compact}{compactos} es
compacta y que una intersección
arbitraria también lo es. (b) Si $K$ es
\omterm{def:b3:topology:compact}{compacto}, $F$ es cerrado y
$K \cap F = \varnothing$, demostrar que
$d(K, F) = \inf\{d(x,y) : x \in K, y \in F\} > 0$; dar un contraejemplo
con dos cerrados disjuntos. (c) Demostrar que $X$ es
\omterm{def:b3:topology:compact}{compacto} si y solo si \emph{toda} $f
\colon X \to \R$ \omterm{def:b3:topology:continuity}{continua} está
acotada. \emph{(Si alguna sucesión no tiene subsucesión convergente,
constrúyase una función \omterm{def:b3:topology:continuity}{continua} no
acotada con soporte cerca de sus términos; o úsense funciones de tipo
$x \mapsto
d(x, \cdot)$ --- una vía limpia: si $(x_n)$ no tiene punto de
acumulación, el conjunto $\{x_n\}$ es cerrado y discreto, y $f(x_n) = n$
se extiende \omterm{def:b3:topology:continuity}{continuamente} sin
necesidad de Tietze:
$f(x) = \sum_n n\,\max\bigl(0, 1 - \frac{d(x, x_n)}{r_n}\bigr)$
para $r_n$ suficientemente pequeño.)}
\end{exercise}

\begin{exercise}[$\star\star$]\label{exo:b3:topology:7}
(Número de Lebesgue) Sea $(U_i)$ un recubrimiento \omterm{def:b3:topology:topology}{abierto} de un espacio
métrico \omterm{def:b3:topology:compact}{compacto} $X$. Demostrar que
existe $\delta > 0$ tal que todo subconjunto de diámetro $< \delta$ está
contenido en un único $U_i$. \emph{(En caso contrario, tómense $A_n$ de
diámetro $< 1/n$ contenidos en ningún $U_i$, y un punto de acumulación
de ciertos $x_n \in A_n$.)} Deducir de nuevo el teorema de Heine
(el~\cref{cor:b3:topology:heineborel}(b)).
\end{exercise}

\begin{exercise}[$\star\star$]\label{exo:b3:topology:8}
(a) Demostrar que $GL_n(\R)$ es un subconjunto \omterm{def:b3:topology:topology}{abierto} y
\omterm{def:b3:topology:interior}{denso} de $M_n(\R)$, y que es
disconexo: el signo del determinante lo separa en (al menos) dos piezas.
Demostrar, en cambio, que $GL_n(\C)$ es
\omterm{def:b3:topology:pathconnected}{conexo por caminos}. \emph{(Para
$A, B \in GL_n(\C)$: $\lambda \mapsto \det\bigl((1 -
\lambda)A + \lambda B\bigr)$ es un polinomio no idénticamente nulo, de
modo que tiene un número finito de raíces en $\C$: tómese un
\omterm{def:b3:topology:pathconnected}{camino} de matrices $\lambda$ en
$\C$ de $0$ a $1$ que las evite.)} (b) Demostrar la densidad:
$A + \varepsilon I$ es invertible para todo
$\varepsilon > 0$.
\end{exercise}

\begin{exercise}[$\star\star$]\label{exo:b3:topology:9}
(a) Demostrar que las \omterm{def:b3:topology:components}{componentes
conexas} de $\Q$ se reducen a puntos y que $\Q$ no es
\omterm{def:b3:topology:locallycompact}{localmente compacto} (un
\omterm{def:b3:topology:topology}{entorno}
\omterm{def:b3:topology:compact}{compacto} de $0$ en $\Q$ contendría
$[-r, r]\cap\Q$, cuya \omterm{def:b3:topology:interior}{clausura} en
$\Q$ no es compacta: córtese en un
irracional). (b) Completar el~\cref{ex:b3:topology:sinecurve}: ningún
\omterm{def:b3:topology:pathconnected}{camino} une $(0,0)$ con $\Gamma$.
\emph{(Si $\gamma = (\gamma_1, \gamma_2)$ es un tal
\omterm{def:b3:topology:pathconnected}{camino}, sea
$t^* = \sup\{t : \gamma_1(t) = 0\}$; para $t
> t^*$ el punto se mueve sobre la gráfica; elíjanse $t_k \downarrow
t^*$ con $\gamma_1(t_k) \to 0$ alcanzando abscisas donde $\sin$ vale
alternativamente $\pm1$ --- el teorema del valor intermedio las
proporciona ---, en contra de la
\omterm{def:b3:topology:continuity}{continuidad} de $\gamma_2$ en
$t^*$.)}
\end{exercise}

\begin{exercise}[$\star\star\star$]\label{exo:b3:topology:10}
El conjunto de Cantor de los tercios centrales $C = \bigcap_n C_n$,
donde $C_0 =
[0,1]$ y $C_{n+1}$ suprime el tercio central \omterm{def:b3:topology:topology}{abierto} de cada intervalo
de $C_n$. (a) Demostrar que
$C = \bigl\{\sum_{n\geq1} a_n3^{-n} : a_n \in
\{0,2\}\bigr\}$, que $C$ es \omterm{def:b3:topology:compact}{compacto}
con \omterm{def:b3:topology:interior}{interior} vacío y que no tiene
puntos aislados (es \emph{\omterm{prop:b3:galois:perfect}{perfecto}}).
(b) Demostrar que $C$ es totalmente disconexo. (c) Demostrar que
$x \mapsto (a_n(x))_n$ es un
\omterm{def:b3:topology:continuity}{homeomorfismo} $C \to
\{0,2\}^\N$ (con la \omterm{def:b3:topology:constructions}{topología
producto} sobre el espacio discreto de dos puntos); deducir que $C$ es
no numerable y que $C \times C
\cong C$.
\end{exercise}

\begin{exercise}[$\star\star\star$]\label{exo:b3:topology:11}
Proyección estereográfica: desde el polo norte $N = (0, \dots,
0, 1)$ de $S^n \subseteq \R^{n+1}$, la aplicación $\sigma(x) =
\frac{(x_1, \dots, x_n)}{1 - x_{n+1}}$ es un
\omterm{def:b3:topology:continuity}{homeomorfismo} $S^n
\setminus \{N\} \to \R^n$ (dese la inversa explícitamente). Deducir que
$S^n$ es \omterm{def:b3:topology:continuity}{homeomorfa} a la
\omterm{def:b3:topology:locallycompact}{compactificación por un punto}
$\widehat{\R^n}$ y que, suprimiendo \emph{cualquier} punto de $S^n$,
queda un espacio \omterm{def:b3:topology:continuity}{homeomorfo} a
$\R^n$.
\end{exercise}

\begin{exercise}[$\star\star\star$]\label{exo:b3:topology:12}
(La \omterm{ex:b3:topology:sinecurve}{curva seno del topólogo}) Sea
\[
T = \bigl\{(x, \sin\tfrac1x) : 0 < x \leq 1\bigr\}
\cup \bigl(\{0\}\times\intcc{-1}1\bigr) \subseteq \R^2 .
\]
(a) Demostrar que $T$ es \omterm{def:b3:topology:compact}{compacto} y
que es la \omterm{def:b3:topology:interior}{clausura} de la parte que es
gráfica. (b) Demostrar que $T$ es
\omterm{def:b3:topology:connected}{conexo} \emph{(la gráfica es
\omterm{def:b3:topology:connected}{conexa} por ser imagen
\omterm{def:b3:topology:continuity}{continua}; su
\omterm{def:b3:topology:interior}{clausura} sigue siendo
\omterm{def:b3:topology:connected}{conexa})}. (c) Demostrar que $T$
\emph{no} es \omterm{def:b3:topology:pathconnected}{conexo por caminos}:
ningún \omterm{def:b3:topology:pathconnected}{camino}
\omterm{def:b3:topology:continuity}{continuo} une $(0,0)$ con
$(1, \sin1)$. \emph{(Si $\gamma =
(\gamma_1, \gamma_2)$ es un tal
\omterm{def:b3:topology:pathconnected}{camino}, sea $t^* = \sup\{t :
\gamma_1(t) = 0\}$; justo después de $t^*$, $\gamma_1$ toma todos los
valores positivos pequeños (teorema del valor intermedio), de modo que
$\gamma_2$ oscila entre $\pm1$ en todo intervalo $(t^*, t^* + \delta)$
--- en contra de la \omterm{def:b3:topology:continuity}{continuidad} en
$t^*$.)} (d) Concluir que la \omterm{def:b3:topology:connected}{conexión}
por \omterm{def:b3:topology:pathconnected}{caminos} es estrictamente más fuerte que la
\omterm{def:b3:topology:connected}{conexión}, y demostrar que ningún
ejemplo de este tipo puede ser \omterm{def:b3:topology:topology}{abierto} en $\R^n$: un subconjunto
\emph{\omterm{def:b3:topology:topology}{abierto}} y \omterm{def:b3:topology:connected}{conexo} de $\R^n$ es
\omterm{def:b3:topology:pathconnected}{conexo por caminos} \emph{(el
conjunto de los puntos unibles al punto \omterm{def:b3:topology:basis}{base} por un
\omterm{def:b3:topology:pathconnected}{camino} es \omterm{def:b3:topology:topology}{abierto} y cerrado en
el dominio)}.
\end{exercise}

\section{Problema: el conjunto de Cantor y una curva que llena el
cuadrado}

\begin{problem}[{Problema de fin de semana --- las curvas de Peano
existen, y por qué
no son homeomorfismos}]\label{pb:b3:topology:1} En 1890 Peano dejó
atónito al análisis con una \emph{sobreyección
\omterm{def:b3:topology:continuity}{continua}} $[0,1] \to [0,1]^2$: una
curva que llena un cuadrado. Construiremos una a mano a partir del
conjunto de Cantor $C$ del~\cref{exo:b3:topology:10} (cuyos resultados
pueden usarse libremente) y demostraremos después que ninguna aplicación
así puede ser inyectiva: los cuadrados no son curvas. En todo el
problema, los elementos de $C$ se escriben $x =
\sum_{n \geq 1} \frac{2b_n(x)}{3^n}$ con dígitos $b_n(x) \in
\{0, 1\}$.

\medskip
\noindent\textbf{Parte I --- Leer los dígitos de manera
\omterm{def:b3:topology:continuity}{continua}.}
\begin{enumerate}
  \item Demostrar que si $x, y \in C$ cumplen $\abs{x - y} <
        3^{-m}$, entonces $b_n(x) = b_n(y)$ para todo $n \leq m$.
  \emph{(Si el primer dígito en que difieren es el $b_k$, entonces
        $\abs{x - y} \geq \frac{2}{3^k} - \sum_{j>k}\frac2{3^j}
        = \frac1{3^k}$.)}
  \item Deducir que cada función dígito $b_n \colon C \to \{0,
        1\}$ es \omterm{def:b3:topology:continuity}{continua} y volver a
  obtener el \omterm{def:b3:topology:continuity}{homeomorfismo} $C
        \cong \{0,1\}^\N$ del~\cref{exo:b3:topology:10}(c).
\end{enumerate}

\medskip
\noindent\textbf{Parte II --- Una sobreyección
\omterm{def:b3:topology:continuity}{continua} $C \to
[0,1]^2$.}
\begin{enumerate}[resume]
  \item Defínase $g \colon C \to [0,1]$ mediante $g(x) = \sum_{n\geq1}
        b_n(x)\,2^{-n}$ (léanse los dígitos de Cantor como binarios).
  Demostrar que $g$ es \omterm{def:b3:topology:continuity}{continua}
  (úsese la pregunta 1) y sobreyectiva. ¿Es inyectiva? \item Defínase
  $\Phi \colon C \to [0,1]^2$ mediante
        \[
        \Phi(x) = \Bigl(\ \sum_{n\geq1} b_{2n-1}(x)\,2^{-n},\
        \sum_{n\geq1} b_{2n}(x)\,2^{-n}\ \Bigr)
        \]
        (los dígitos impares dan la abscisa y los pares la ordenada).
        Demostrar que $\Phi$ es
        \omterm{def:b3:topology:continuity}{continua} y sobreyectiva.
\end{enumerate}

\medskip
\noindent\textbf{Parte III --- Rellenar los huecos: la curva de Peano.}
\begin{enumerate}[resume]
  \item El complementario $[0,1] \setminus C$ es una unión numerable de
  intervalos \omterm{def:b3:topology:topology}{abiertos} disjuntos $(u, v)$ (los tercios centrales
  suprimidos) cuyos extremos están en $C$. Defínase $f \colon
        [0,1] \to [0,1]^2$ como $f = \Phi$ sobre $C$, extendida
  \emph{afínmente} en cada hueco:
        \[
        f(t) = \frac{v - t}{v - u}\,\Phi(u) +
        \frac{t - u}{v - u}\,\Phi(v),
        \qquad t \in (u, v).
        \]
        Demostrar que $f$ está bien definida y es sobreyectiva sobre
        $[0,1]^2$.
  \item Demostrar que $f$ es
  \omterm{def:b3:topology:continuity}{continua} en todo punto de $[0,1]
        \setminus C$ (es localmente afín) y en todo punto de $C$: dado
  $\varepsilon$, tómese $m$ con $2^{-m} <
        \varepsilon/2$ y úsese la pregunta 1 para controlar $\Phi$ sobre
  $C$ cerca de $x$; compruébese después que la interpolación afín no
  puede escaparse: sobre un hueco $(u,v) \subseteq
        (x - 3^{-m}, x + 3^{-m})$, los valores $f(t)$ están en el
  segmento $[\Phi(u), \Phi(v)]$, cuyos dos extremos están próximos a
  $\Phi(x)$. Concluir: $f$ es una \emph{sobreyección
  \omterm{def:b3:topology:continuity}{continua}} $[0,1] \to [0,1]^2$.
  \item Deducir la existencia de sobreyecciones
  \omterm{def:b3:topology:continuity}{continuas} $[0,1] \to [0,1]^n$
  para todo $n \geq 2$, y $\R \to \R^2$.
\end{enumerate}

\medskip
\noindent\textbf{Parte IV --- Pero nunca inyectiva.}
\begin{enumerate}[resume]
  \item Demostrar que una \emph{inyección}
  \omterm{def:b3:topology:continuity}{continua} $f \colon [0,1]
        \to [0,1]^2$ sería un
  \omterm{def:b3:topology:continuity}{homeomorfismo} sobre su imagen
        (\cref{cor:b3:topology:compacthomeo}).
  \item Demostrar que $[0,1]$ y $[0,1]^2$ no son
  \omterm{def:b3:topology:continuity}{homeomorfos}: suprímase un punto
  bien elegido y compárese la
  \omterm{def:b3:topology:connected}{conexión}
        (\cref{met:b3:topology:connectargument}).
  \item Concluir: una sobreyección
  \omterm{def:b3:topology:continuity}{continua} $[0,1] \to [0,1]^2$
  nunca puede ser inyectiva --- si lo fuera, la pregunta 8 haría
  $[0,1]^2 = f([0,1])$ \omterm{def:b3:topology:continuity}{homeomorfo} a
  $[0,1]$, en contra de la pregunta 9. ¿Dónde intervino exactamente la
  \omterm{def:b3:topology:compact}{compacidad} de $[0,1]$ en el argumento? \item (Culminación) Reúnase la
  moraleja: existe una sobreyección
  \omterm{def:b3:topology:continuity}{continua}, pero ninguna biyección
  \omterm{def:b3:topology:continuity}{continua} $[0,1] \to
        [0,1]^2$. ¿Qué dice esto sobre la «dimensión» como noción
  \omterm{def:b3:topology:topology}{topológica}? Formúlese con precisión
  un teorema demostrado en este problema y un enunciado plausible que
  queda fuera de nuestras herramientas (la invariancia del dominio).
\end{enumerate}

\medskip
\noindent\textbf{Parte V --- La aritmética del conjunto de Cantor.}
\begin{enumerate}[resume]
  \item Demostrar que $C + C = [0, 2]$: dado $t \in [0, 1]$, escríbase
  $t = \sum_n d_n3^{-n}$ con dígitos $d_n \in \{0,
        1, 2\}$ y sepárese cada dígito como $d_n = a_n + b_n$ con
  $a_n, b_n \in \{0, 1\}$; concluir que $\frac C2 +
        \frac C2 = [0, 1]$ \emph{(¿qué subconjunto de $[0,1]$ es
  $\frac C2$, en términos de dígitos ternarios?)} y reescálese después.
  Nunca hacen falta acarreos --- explíquese por qué ese es el meollo.
  \item Interprétese geométricamente: el «polvo de Cantor» $C \times
        C \subseteq \R^2$, de
  \omterm{def:b3:topology:interior}{interior} vacío, proyecta una sombra
  \emph{completa} sobre la diagonal: la proyección $(x,
        y) \mapsto x + y$ lo aplica sobre $[0, 2]$. Regístrese también
  la autosemejanza $C = \frac C3 \cup \bigl(\frac C3 +
        \frac23\bigr)$, la ecuación que hay detrás de toda imagen de
        $C$.
  \item Demostrar que el entrelazado de dígitos define un
  \omterm{def:b3:topology:continuity}{homeomorfismo} $C
        \to C \times C$ y deducir $C \cong C^n$ para todo $n$ --- el
  conjunto de Cantor es su propio cuadrado, su propio cubo, \dots{} ¿Qué
  espacios familiares comparten esta propiedad? \item Demostrar que
  $\frac14 \in C$ \emph{(calcúlese su desarrollo ternario:
  $\frac14 = \sum_{k\geq1} 2\cdot3^{-2k}$)} aunque $\frac14$ no sea
  extremo de ningún intervalo suprimido; deducir --- contando extremos
  --- que los extremos forman un subconjunto numerable y propio de $C$.
  \item Demostrar que la aplicación de lectura binaria $g$ de la
  pregunta 3 es a lo sumo $2$ a $1$, y describir exactamente qué puntos
  de $[0,1]$ tienen dos preimágenes. ($g$ colapsa $C$ sobre $[0,1]$
  pegando una infinidad numerable de parejas: la sombra combinatoria de
  la escalera de Cantor, que reaparecerá en el~\cref{ch:b3:measure}.)
\end{enumerate}

\medskip
\noindent\textbf{Parte VI --- Conjuntos
\omterm{def:b3:topology:compact}{compactos}
\omterm{prop:b3:galois:perfect}{perfectos}: Cantor en todas partes.} Un
espacio métrico \omterm{def:b3:topology:compact}{compacto} no vacío $P$
es \emph{\omterm{prop:b3:galois:perfect}{perfecto}} si no tiene puntos
aislados.
\begin{enumerate}[resume]
  \item Sea $P$ \omterm{def:b3:topology:compact}{compacto}
  \omterm{prop:b3:galois:perfect}{perfecto} y sean $x \in P$, $r > 0$.
  Demostrar que la bola $B(x, r)$ contiene dos puntos $y_0
        \neq y_1$ de $P$ y, por tanto, dos bolas cerradas
  \emph{disjuntas} $\bar B(y_0, \rho)$, $\bar B(y_1, \rho)$ dentro de
  $B(x, r)$, centradas en puntos de $P$ y de radio $\rho
        \leq r/4$ tan pequeño como se quiera. Explíquese por qué la
  perfección (la ausencia de puntos aislados) es exactamente lo que
  permite repetir esta división dentro de \emph{cada} una de las dos
  bolas nuevas. \item Itérese: constrúyanse cerrados
  $F_w \neq \varnothing$ indexados por palabras binarias finitas $w$,
  con $F_{w0},
        F_{w1} \subseteq F_w$ disjuntos y
  $\operatorname{diam}F_w \leq 2^{-\abs w}
        \operatorname{diam}P$. Demostrar que, para toda palabra infinita
  $\varepsilon \in \{0,1\}^\N$, la intersección
  $\bigcap_mF_{\varepsilon_1\cdots\varepsilon_m}$ es un único punto
  $h(\varepsilon)$ \emph{(propiedad de intersección finita de los
  \omterm{def:b3:topology:compact}{compactos} $P$)}. \item Demostrar que
  $h\colon\{0,1\}^\N \to P$ es inyectiva y
  \omterm{def:b3:topology:continuity}{continua}, y concluir: \emph{todo
  espacio métrico \omterm{def:b3:topology:compact}{compacto}
  \omterm{prop:b3:galois:perfect}{perfecto} es no numerable} --- de
  hecho, de cardinal al menos el de $\{0,1\}^\N$. Recupérese: $[0,1]$ y
  $C$ son no numerables. \item Demostrar que $h$ es un
  \omterm{def:b3:topology:continuity}{homeomorfismo} sobre su imagen
  \emph{(inyección \omterm{def:b3:topology:continuity}{continua} desde
  un \omterm{def:b3:topology:compact}{compacto},
  el~\cref{cor:b3:topology:compacthomeo})}: todo espacio métrico
  \omterm{def:b3:topology:compact}{compacto}
  \omterm{prop:b3:galois:perfect}{perfecto} contiene una copia
  \omterm{def:b3:topology:continuity}{homeomorfa} del conjunto de
  Cantor. El conjunto de Cantor no es una rareza, sino el germen
  universal de la perfección compacta.
  \item Deducir que todo espacio métrico
  \omterm{def:b3:topology:compact}{compacto} numerable tiene un punto
  aislado, y exhibir uno en el que los puntos aislados sean
  \omterm{def:b3:topology:interior}{densos} sin serlo todo:
  $K = \{0\} \cup
        \{\frac1n : n \geq 1\}$. \item (Cantor--Bendixson en $\R$) Sea
  $F \subseteq \R$ cerrado. Diremos que $x$ es un \emph{punto de
  condensación} de $F$ si todo
  \omterm{def:b3:topology:topology}{entorno} de $x$ corta a $F$ en un
  conjunto no numerable. Demostrar que los puntos de condensación de un
  cerrado no numerable $F$ forman un cerrado
  \omterm{prop:b3:galois:perfect}{perfecto} no vacío $P$ y que $F
        \setminus P$ es numerable \emph{(recúbranse los puntos que no
  son de condensación mediante una cantidad numerable de intervalos
  racionales que cortan a $F$ en conjuntos numerables)}. Concluir: todo
  subconjunto cerrado de $\R$ es numerable o tiene el cardinal del
  \omterm{def:b3:topology:continuity}{continuo} --- la hipótesis del \omterm{def:b3:topology:continuity}{continuo} se cumple para los cerrados.
\end{enumerate}

\medskip
\noindent\textbf{Parte VII --- Codas: cuán regular, y cuán lejos de ser
\omterm{prop:b3:galois:perfect}{perfecto}.}
\begin{enumerate}[resume]
  \item (Regularidad Hölder de la curva) Póngase $\alpha =
        \frac{\ln 2}{2\ln 3} \approx 0.3155$ y obsérvese la identidad
  $3^\alpha = \sqrt2$. Usando la pregunta 1, demostrar que
        \[
        \norm{\Phi(x) - \Phi(y)}_\infty \leq 2\,\abs{x -
        y}^\alpha \qquad (x, y \in C),
        \]
        propáguese después la estimación por los huecos afines:
        demuéstrese que la curva de Peano $f$ de la pregunta 6 cumple
        $\norm{f(t) - f(s)}_\infty \leq 6\abs{t - s}^\alpha$ en todo
        $[0,1]$ \emph{(trátese un par situado en un mismo hueco por
        interpolación, y un par general pasando por los puntos extremos
        de $C \cap [t, s]$)}. \item (El exponente $\frac12$ es un muro)
        Demostrar que ninguna sobreyección $F \colon [0,1] \to [0,1]^2$
        puede ser Hölder de exponente $\alpha$ con $\alpha > \frac12$:
        córtese $[0,1]$ en $N$ intervalos, acótense los diámetros de sus
        imágenes y cuéntense los puntos de la retícula
        $\bigl(\frac iM, \frac jM\bigr)$, $0 \leq i, j \leq M$,
        que puede contener un conjunto de diámetro $< \frac1M$; tómese
        $M$ del orden de $N^\alpha$ y hágase $N \to \infty$. Sitúese
        nuestra curva ($\alpha \approx 0.3155$) respecto del muro y
        regístrese sin demostración que la curva de Hilbert alcanza el
        exponente crítico $\frac12$. \item (Conjuntos derivados: medir
        la imperfección) Para $F$ cerrado en $\R$, sea $F'$ el conjunto
        de sus puntos de acumulación $F$ (un cerrado), e itérese:
        $F^{(0)} = F$, $F^{(i+1)}
        = (F^{(i)})'$. Comprobar que
        \[
        K = \{0\} \cup \Bigl\{\frac1m\Bigr\}_{m \geq 1} \cup
        \Bigl\{\frac1m + \frac1n : m \geq 1,\ n > m^2\Bigr\}
        \]
        es un \omterm{def:b3:topology:compact}{compacto} numerable con
        $K' = \{0\} \cup
        \{\frac1m\}$, $K'' = \{0\}$, $K''' = \varnothing$
        \emph{(compruébese que el $m$-ésimo racimo vive en el intervalo
        $(\frac1m, \frac1{m-1})$, de modo que los racimos no se
        entrelazan)} y que sus puntos aislados son
        \omterm{def:b3:topology:interior}{densos} en $K$, como predice
        la pregunta 21. Descríbase la inducción que produce, para todo
        $j$, un \omterm{def:b3:topology:compact}{compacto} numerable
        $K_j$ con $K_j^{(j)} = \{0\}$ y $K_j^{(j+1)} =
        \varnothing$, y contrástese con los conjuntos
        \omterm{prop:b3:galois:perfect}{perfectos} de la pregunta 22,
        para los cuales la derivación nunca se mueve: el rango finito es
        exactamente lo contrario de la perfección.
\end{enumerate}
\end{problem}
